Transcranial focused ultrasound stimulation (tFUS) can modulate human brain activity, but its impact on large-scale functional connectivity (FC) \emph{in vivo} remains incompletely characterized. Here we asked whether brief tFUS to the subgenual anterior cingulate cortex (sgACC) modulates its FC with other brain regions. In a sham-controlled, within-subject study of healthy adults ($N=16$), we measured resting-state blood-oxygen-level dependent (BOLD) activity with functional magnetic resonance imaging (fMRI) before, during, and after five minutes of tFUS (i.e., five 20-second sonications followed by 40 second pauses). Subject-level sgACC–whole-brain FC was analyzed with a linear mixed-effects model (fixed: time, condition, time×condition; random: subject). tFUS increased sgACC-whole brain connectivity after sonication relative to sham. A smaller, non-significant trend was observed during sonication. Baseline-controlled analyses were employed to rule out the possibility that the observed increase in FC was due to a ``regression to the mean'' effect. Interestingly, network-level analysis revealed a clear disassociation in the evolution of resting-state FC across the 15 minute recordings: during sham sessions, the sgACC's connectivity with the default mode network (DMN) increased, while FC with other networks was largely unchanged. On the other hand, active tFUS led to a notable increase in FC with non-DMN networks (especially the cognitive control network), while connectivity with the DMN was stable. These results provide whole-brain evidence that tFUS reconfigures network dynamics in the human brain, motivating clinical applications to disorders characterized by dysregulated brain activity.

